of @@M0@@ in @@M1@@, assumed everywhere to have rank @@M2@@. In loc. cit., this formula follows from a very simple computation of @@M3@@ sheaves. By what we saw in 4.2, the preceding formula implies the same formula in @@M4@@, but only \uline{modulo torsion}. It is quite extraordinary that formula (4.6) in @@M5@@ itself has not been proved at present\SGAAuthorFootnote{*}{cf. the note on page 23.\newline (**)}, nor even the special case obtained by taking @@M6@@, which gives the self-intersection of @@M7@@:
\sgaStableTarget{sga6.e14.eq.4-7}@@M8@@
even when @@M9@@ is the field of complex numbers!

An analogous problem arises when one attempts to compute the Chow ring of the variety @@M10@@ obtained by blowing up @@M11@@ along @@M12@@, imitating the method used to compute @@M13@@ [Exp. VII 4.8], which also works (without tensoring with @@M14@@) for the computation of @@M15@@. The key formula for this determination (with the notation of loc. cit.) is
\sgaStableTarget{sga6.e14.eq.4-8}@@M16@@
a formula currently established in @@M17@@ (for @@M18@@) by a computation of @@M19@@ (Exp. VII 3.4), whose validity in @@M20@@ (for @@M21@@) remains to be established.

Observe that even the @@M22@@-adic analogues of formulas (4.6), (4.7), (4.8) have not been proved at present; when @@M23@@ is the cohomology class of an algebraic cycle, (4.6) and (4.8) hold in any case mod. torsion.

\sgaStableTarget{sga6.e14.para.4-5}\textbf{4.5.} In this section and the following ones, we shall give examples in which, in diagram (4.1), the homomorphism @@M24@@, respectively the homomorphism @@M25@@, is not an isomorphism. For a given @@M26@@, it clearly amounts to the same thing to say that
@@M27@@
is an isomorphism, or a monomorphism, or an epimorphism. Moreover, the image of @@M28@@ is the subring of @@M29@@ generated by the Chern classes of locally free sheaves on @@M30@@. When @@M31@@ is generated as an algebra over @@M32@@ by the classes of invertible sheaves on @@M33@@, one sees immediately that the subring of @@M34@@ under consideration is also the subring generated by @@M35@@, the component of degree @@M36@@ of @@M37@@.

